\textbf{Variáveis e suas dimensões:}

\begin{center}
\begin{tabular}{lll}
\hline
Variável & Símbolo & Dimensões \\
\hline
Parâmetro catenário & $c$ & $[L]$ \\
Tensão horizontal & $T_0$ & $[M\,L\,T^{-2}]$ (força) \\
Peso por unidade de comprimento & $\gamma$ & $[M\,L^{-1}\,T^{-2}]$ \\
\hline
\end{tabular}
\end{center}

\textbf{Aplicação do Teorema de Buckingham--$\Pi$:}

São $n=3$ variáveis e $k=3$ dimensões ($M,L,T$), mas apenas $n-k=3-3=0$ grupos adimensionais independentes. Isso indica que a relação entre as variáveis é \emph{única} (a menos de uma constante adimensional).

\textbf{Dedução por inspecção dimensional:}

Buscamos $c=K\,T_0^a\,\gamma^b$ com dimensões $[L]$:
\[
[L] = [M\,L\,T^{-2}]^a\,[M\,L^{-1}\,T^{-2}]^b.
\]

Sistema:
\[
M:\;a+b=0,\quad T:\;-2a-2b=0\;(\text{redundante}),\quad L:\;a-b=1.
\]

De $a+b=0$: $b=-a$. Substituindo em $a-b=1$: $a-(-a)=2a=1$, logo $a=1/2$ e $b=-1/2$.

Portanto:
\[
c\propto\frac{T_0^{1/2}}{\gamma^{1/2}}=\sqrt{\frac{T_0}{\gamma}}.
\]

\textbf{Verificação a partir da equação da catenária:}

A equação diferencial de equilíbrio do cabo é:
\[
\frac{d^2y}{dx^2}=\frac{\gamma}{T_0}\sqrt{1+\left(\frac{dy}{dx}\right)^2}.
\]

A solução é $y=c\cosh(x/c)$ com $c=T_0/\gamma$. Confirmando:
\[
c = \frac{T_0}{\gamma}.\quad\checkmark
\]

\textbf{Verificação dimensional:}
\[
\left[\frac{T_0}{\gamma}\right]=\frac{M\,L\,T^{-2}}{M\,L^{-1}\,T^{-2}}=L.\quad\checkmark
\]